% ===========================================================================
\chapter{Mapas de Empatia das Personas}
\label{cap_mapas_empatia}
% ===========================================================================

O Mapa de Empatia consolida o mergulho psicológico na experiência subjetiva de cada persona, permitindo que a equipe de engenharia compreenda não apenas as ações operacionais desempenhadas, mas também as tensões sensoriais, expectativas implícitas e emoções subjacentes à interação com a plataforma.

Neste capítulo, são apresentados os mapas de empatia em seis quadrantes estruturados para as três personas do sistema: a discente Beatriz Matos (Seção \ref{sec_mapa_beatriz}), o docente Prof. Dr. Ricardo Sobral (Seção \ref{sec_mapa_ricardo}) e o administrador de TI Carlos Eduardo Santos (Seção \ref{sec_mapa_carlos}), culminando em uma análise comparativa dos perfis (Seção \ref{sec_sintese_mapas}).

% ===========================================================================
\section{Mapa de Empatia da Persona Discente: Beatriz Matos}
\label{sec_mapa_beatriz}
% ===========================================================================

A Figura \ref{fig_mapa_empatia_beatriz} ilustra o Mapa de Empatia estruturado para a persona primária do SIGEA-GTT.

\begin{figure}[!ht]
\centering
\resizebox{\textwidth}{!}{%
\begin{tikzpicture}[
    box/.style={draw=clinical-700, fill=white, rounded corners=5pt, line width=1pt, inner sep=8pt, text width=6.8cm, font=\scriptsize, align=left},
    painbox/.style={draw=danger-red, fill=danger-red!10, rounded corners=5pt, line width=1pt, inner sep=8pt, text width=6.8cm, font=\scriptsize, align=left},
    gainbox/.style={draw=hospital-green, fill=hospital-green!10, rounded corners=5pt, line width=1pt, inner sep=8pt, text width=6.8cm, font=\scriptsize, align=left},
    header/.style={fill=clinical-700, text=white, font=\bfseries\footnotesize, rounded corners=3pt, inner sep=4pt, align=center},
    centerpic/.style={draw=clinical-700, fill=clinical-100!40, circle, line width=2pt, inner sep=6pt, align=center, font=\bfseries\small}
]

% Persona Central
\node[centerpic] (center) at (0,0) {
    \scalebox{2.0}{\color{clinical-700}$\mathbf{\Omega}$} \\
    \vspace{2pt}
    BEATRIZ MATOS \\
    \tiny Discente de Enfermagem
};

% Quadrante 1: O que Pensa e Sente? (Norte)
\node[box, above=1.2cm of center] (pensa) {
    \textbf{Preocupações e Aspirações}:
    \begin{itemize}\setlength{\itemsep}{1pt}
        \item ``Será que conseguirei auditar o prontuário em menos de 20 minutos sem me desesperar?''
        \item ``Não quero errar a classificação de severidade do dano na escala NCC MERP.''
        \item Deseja sentir segurança para atuar nos plantões do internato hospitalar;
        \item Preocupa-se em entender as falhas nos processos em vez de culpar colegas.
    \end{itemize}
};
\node[header, above=0.1cm of pensa] {O QUE PENSA E SENTE?};

% Quadrante 2: O que Escuta? (Oeste)
\node[box, left=0.8cm of center] (escuta) {
    \textbf{Vozes do Ambiente}:
    \begin{itemize}\setlength{\itemsep}{1pt}
        \item Professores afirmando que auditoria clínica exige agilidade mental e olhar treinado;
        \item Colegas de estágio comentando o medo de serem advertidos em caso de erro;
        \item Enfermeiros supervisores cobrando atenção redobrada em prescrições complexas;
        \item Diretrizes do COFEN e da OMS sobre metas internacionais de segurança.
    \end{itemize}
};
\node[header, above=0.1cm of escuta] {O QUE ESCUTA?};

% Quadrante 3: O que Vê? (Leste)
\node[box, right=0.8cm of center] (ve) {
    \textbf{Ambiente e Estímulos Visuais}:
    \begin{itemize}\setlength{\itemsep}{1pt}
        \item Prontuários simulados repletos de laudos, prescrições e notas de enfermagem;
        \item Monitores de computador nos laboratórios da UFS com textos densos;
        \item Cronômetro do celular correndo enquanto tenta encontrar evidências textuais;
        \item Fórmulas complexas e tabelas dispersas de Ishikawa e Matriz GUT.
    \end{itemize}
};
\node[header, above=0.1cm of ve] {O QUE VÊ?};

% Quadrante 4: O que Fala e Faz? (Sul)
\node[box, below=1.2cm of center] (fala) {
    \textbf{Comportamento Observável}:
    \begin{itemize}\setlength{\itemsep}{1pt}
        \item Discute casos clínicos em grupo usando terminologia padronizada de enfermagem;
        \item Destaca trechos do prontuário com marca-texto e copia trechos em rascunhos;
        \item Tenta justificar a relação de causa e efeito entre medicações e sintomas;
        \item Pede esclarecimentos ao professor sobre falsos positivos em gatilhos laboratoriais.
    \end{itemize}
};
\node[header, above=0.1cm of fala] {O QUE FALA E FAZ?};

% Pain Box: Dores (Inferior Esquerdo)
\node[painbox, below=0.6cm of fala.south west, anchor=north west] (dores) {
    \textbf{\color{danger-red} DORES (FRUSTRAÇÕES E OBSTÁCULOS)}:
    \begin{itemize}\setlength{\itemsep}{1pt}
        \item Ansiedade aguda gerada por contadores de tempo abruptos sem avisos graduais;
        \item Sobrecarga de leitura em prontuários desorganizados sem abas tabuladas;
        \item Medo de retaliação e cultura punitiva tradicional hospitalar;
        \item Dificuldade em memorizar os 53 gatilhos canônicos do IHI sem um catálogo ágil.
    \end{itemize}
};

% Gain Box: Ganhos (Inferior Direito)
\node[gainbox, below=0.6cm of fala.south east, anchor=north east] (ganhos) {
    \textbf{\color{hospital-green} GANHOS (DESEJOS E ASPIRAÇÕES)}:
    \begin{itemize}\setlength{\itemsep}{1pt}
        \item Concluir a auditoria dentro dos 20 minutos com clareza e sem pânico;
        \item Interface visual com alertas suaves de tempo aos 15 e 18 minutos;
        \item Prontuário organizado em abas com vinculação direta de evidências citadas;
        \item Acesso rápido ao Catálogo Didático de Gatilhos com busca textual (\textit{debounce});
        \item Receber nota justa e parecer pedagógico formativo detalhado do professor titular.
    \end{itemize}
};

\end{tikzpicture}
}
\caption{Mapa de Empatia em Seis Quadrantes: Beatriz Matos (Discente)}
\label{fig_mapa_empatia_beatriz}
\fonte{Elaborado pelo autor (2026).}
\end{figure}

\subsection{Análise do Universo Psicológico e Sensorial de Beatriz}
O mapa de empatia revela que o principal vetor de estresse de Beatriz decorre da dicotomia entre a necessidade de rigor analítico e a pressão temporal estrita dos 20 minutos preconizada pelo IHI. Quando submetida a cronômetros convencionais que disparam alarmes ruidosos sem pré-aviso, a discente experimenta um bloqueio de atenção seletiva, deixando passar gatilhos óbvios em laudos laboratoriais (ex.: creatinina sérica elevada ou queda súbita de hemoglobina).

Para mitigar essa dor, a aplicação projeta o prontuário em abas organizadas no componente \texttt{AuditWorkspaceComponent}, acoplando um cronômetro regressivo com indicador visual discreto que transiciona de azul clínico (\texttt{\#2563EB}) para âmbar (\texttt{\#D97706}) aos 15 minutos e carmim cirúrgico (\texttt{\#DC2626}) aos 18 minutos, fornecendo previsibilidade e autocontrole emocional à discente.

% ===========================================================================
\section{Mapa de Empatia da Persona Docente: Prof. Dr. Ricardo Sobral}
\label{sec_mapa_ricardo}
% ===========================================================================

A Figura \ref{fig_mapa_empatia_ricardo} apresenta o Mapa de Empatia estruturado para o docente titular e auditor clínico.

\begin{figure}[!ht]
\centering
\resizebox{\textwidth}{!}{%
\begin{tikzpicture}[
    box/.style={draw=hospital-green, fill=white, rounded corners=5pt, line width=1pt, inner sep=8pt, text width=6.8cm, font=\scriptsize, align=left},
    painbox/.style={draw=danger-red, fill=danger-red!10, rounded corners=5pt, line width=1pt, inner sep=8pt, text width=6.8cm, font=\scriptsize, align=left},
    gainbox/.style={draw=hospital-green, fill=hospital-green!10, rounded corners=5pt, line width=1pt, inner sep=8pt, text width=6.8cm, font=\scriptsize, align=left},
    header/.style={fill=hospital-green, text=white, font=\bfseries\footnotesize, rounded corners=3pt, inner sep=4pt, align=center},
    centerpic/.style={draw=hospital-green, fill=hospital-green!20, circle, line width=2pt, inner sep=6pt, align=center, font=\bfseries\small}
]

% Persona Central
\node[centerpic] (center) at (0,0) {
    \scalebox{2.0}{\color{hospital-green}$\mathbf{\Psi}$} \\
    \vspace{2pt}
    PROF. DR. RICARDO \\
    \tiny Docente e Auditor Sênior
};

% Quadrante 1: O que Pensa e Sente?
\node[box, above=1.2cm of center] (pensa) {
    \textbf{Preocupações e Aspirações}:
    \begin{itemize}\setlength{\itemsep}{1pt}
        \item ``Meus estudantes precisam sair da universidade dominando o raciocínio causal do IHI.''
        \item Sente exaustão profunda ao corrigir pilhas intermináveis de relatórios impressos;
        \item Preocupa-se com a fidedignidade dos cálculos de taxas de eventos adversos;
        \item Quer que cada aluno receba um feedback individualizado e formativo.
    \end{itemize}
};
\node[header, above=0.1cm of pensa] {O QUE PENSA E SENTE?};

% Quadrante 2: O que Escuta?
\node[box, left=0.8cm of center] (escuta) {
    \textbf{Vozes do Ambiente}:
    \begin{itemize}\setlength{\itemsep}{1pt}
        \item A coordenação pedagógica cobrando prazos finais para consolidação de notas semestrais;
        \item Alunos solicitando revisão de critérios de correção de atividades práticas;
        \item Comitês de segurança hospitalar discutindo novas diretrizes e metas de qualidade;
        \item Especialistas do IHI recomendando a padronização de indicadores por 1.000 pacientes-dia.
    \end{itemize}
};
\node[header, above=0.1cm of escuta] {O QUE ESCUTA?};

% Quadrante 3: O que Vê?
\node[box, right=0.8cm of center] (ve) {
    \textbf{Ambiente e Estímulos Visuais}:
    \begin{itemize}\setlength{\itemsep}{1pt}
        \item Dezenas de folhas de papel espalhadas na sua mesa de trabalho no DCOMP/Enfermagem;
        \item Planilhas do Excel com dezenas de abas, macros instáveis e fórmulas quebradas;
        \item Monitores Ultrawide na sua sala que poderiam exibir dados com muito mais riqueza;
        \item Relatórios heterogêneos entregues por estudantes com formatações inconsistentes.
    \end{itemize}
};
\node[header, above=0.1cm of ve] {O QUE VÊ?};

% Quadrante 4: O que Fala e Faz?
\node[box, below=1.2cm of center] (fala) {
    \textbf{Comportamento Observável}:
    \begin{itemize}\setlength{\itemsep}{1pt}
        \item Ministra aulas expositivas e conduz simulações clínicas com foco em evidências;
        \item Elabora casos clínicos fictícios complexos para desafiar o discernimento dos alunos;
        \item Compara as entregas com o gabarito oficial anotando discrepâncias à caneta;
        \item Defende veementemente que a segurança do paciente deve focar em melhorias do sistema.
    \end{itemize}
};
\node[header, above=0.1cm of fala] {O QUE FALA E FAZ?};

% Pain Box: Dores
\node[painbox, below=0.6cm of fala.south west, anchor=north west] (dores) {
    \textbf{\color{danger-red} DORES (FRUSTRAÇÕES E OBSTÁCULOS)}:
    \begin{itemize}\setlength{\itemsep}{1pt}
        \item Perda de dezenas de horas semestrais em conferência mecânica e manual de submissões;
        \item Risco constante de erros de digitação e fórmulas em planilhas eletrônicas;
        \item Turmas criadas com nomes confusos ou despadronizados no sistema da instituição;
        \item Dificuldade de verificar se todos os discentes receberam avaliação antes de encerrar a turma.
    \end{itemize}
};

% Gain Box: Ganhos
\node[gainbox, below=0.6cm of fala.south east, anchor=north east] (ganhos) {
    \textbf{\color{hospital-green} GANHOS (DESEJOS E ASPIRAÇÕES)}:
    \begin{itemize}\setlength{\itemsep}{1pt}
        \item Espelho de correção automatizado comparando a submissão discente com o gabarito docente;
        \item Cálculo automático da Sensibilidade, Especificidade e concordância de severidade;
        \item Dashboard com indicadores oficiais do IHI ($TI_{1000}$, $TI_{100}$, $P_{EA}$) gerados em tempo real;
        \item Bloqueio sistêmico para encerramento de turma quando houver notas pendentes de atribuição;
        \item Possibilidade de emitir parecer formativo individualizado em campo textual dedicado.
    \end{itemize}
};

\end{tikzpicture}
}
\caption{Mapa de Empatia em Seis Quadrantes: Prof. Dr. Ricardo Sobral (Docente)}
\label{fig_mapa_empatia_ricardo}
\fonte{Elaborado pelo autor (2026).}
\end{figure}

\subsection{Análise do Universo Psicológico e Sensorial de Ricardo}
A dor central do Prof. Dr. Ricardo Sobral não é pedagógica, mas operacional: a rotina mecânica de comparar linha a linha as submissões dos estudantes contra o gabarito oficial drena a energia que deveria ser dedicada à mentoria formativa. Além disso, a consolidação manual das equações epidemiológicas do IHI em planilhas do Excel representa um ponto único de falha analítica.

O SIGEA-GTT responde a essas dores no módulo de avaliação docente (\texttt{UC19}, \texttt{UC20} e \texttt{UC21}), fornecendo o espelho de correção automatizado, o cálculo instantâneo de indicadores e a validação invariante que proíbe o encerramento do semestre com pendências avaliativas.

% ===========================================================================
\section{Mapa de Empatia do Administrador: Carlos Eduardo Santos}
\label{sec_mapa_carlos}
% ===========================================================================

A Figura \ref{fig_mapa_empatia_carlos} apresenta o Mapa de Empatia estruturado para o analista de TI do DCOMP.

\begin{figure}[!ht]
\centering
\resizebox{\textwidth}{!}{%
\begin{tikzpicture}[
    box/.style={draw=clinical-500, fill=white, rounded corners=5pt, line width=1pt, inner sep=8pt, text width=6.8cm, font=\scriptsize, align=left},
    painbox/.style={draw=danger-red, fill=danger-red!10, rounded corners=5pt, line width=1pt, inner sep=8pt, text width=6.8cm, font=\scriptsize, align=left},
    gainbox/.style={draw=hospital-green, fill=hospital-green!10, rounded corners=5pt, line width=1pt, inner sep=8pt, text width=6.8cm, font=\scriptsize, align=left},
    header/.style={fill=clinical-500, text=white, font=\bfseries\footnotesize, rounded corners=3pt, inner sep=4pt, align=center},
    centerpic/.style={draw=clinical-500, fill=clinical-500!20, circle, line width=2pt, inner sep=6pt, align=center, font=\bfseries\small}
]

% Persona Central
\node[centerpic] (center) at (0,0) {
    \scalebox{2.0}{\color{clinical-500}$\mathbf{\Delta}$} \\
    \vspace{2pt}
    CARLOS SANTOS \\
    \tiny Analista de TI / DCOMP
};

% Quadrante 1: O que Pensa e Sente?
\node[box, above=1.2cm of center] (pensa) {
    \textbf{Preocupações e Aspirações}:
    \begin{itemize}\setlength{\itemsep}{1pt}
        \item ``O sistema precisa se auto-proteger contra cadastros espúrios e e-mails não institucionais.''
        \item Preocupa-se com a conformidade estrita em relação à LGPD e segurança de dados de alunos;
        \item Quer minimizar chamados de suporte técnico de redefinição de senhas esquecidas;
        \item Deseja interfaces com tabelas limpas, paginação fixa e busca rápida indexada.
    \end{itemize}
};
\node[header, above=0.1cm of pensa] {O QUE PENSA E SENTE?};

% Quadrante 2: O que Escuta?
\node[box, left=0.8cm of center] (escuta) {
    \textbf{Vozes do Ambiente}:
    \begin{itemize}\setlength{\itemsep}{1pt}
        \item Direção do CCET e DCOMP exigindo conformidade de segurança e estabilidade;
        \item Usuários reclamando de sistemas universitários legados lentos e confusos;
        \item Equipes de segurança da informação alertando sobre ataques de enumeração e injeção;
        \item Professores pedindo agilidade no cadastro de turmas e módulos no início do semestre.
    \end{itemize}
};
\node[header, above=0.1cm of escuta] {O QUE ESCUTA?};

% Quadrante 3: O que Vê?
\node[box, right=0.8cm of center] (ve) {
    \textbf{Ambiente e Estímulos Visuais}:
    \begin{itemize}\setlength{\itemsep}{1pt}
        \item Logs de servidores e terminais Linux em estações WUXGA com dados estruturados;
        \item Fila de tickets de atendimento no sistema de suporte acadêmico da UFS;
        \item Banco de dados relacional PostgreSQL com índices e chaves primárias UUID;
        \item Telas de parametrização clínica que precisam ser intuitivas para não médicos.
    \end{itemize}
};
\node[header, above=0.1cm of ve] {O QUE VÊ?};

% Quadrante 4: O que Fala e Faz?
\node[box, below=1.2cm of center] (fala) {
    \textbf{Comportamento Observável}:
    \begin{itemize}\setlength{\itemsep}{1pt}
        \item Aplica boas práticas de segurança, backups diários e monitoramento de desempenho;
        \item Configura parâmetros do sistema com foco em estabilidade e idempotência;
        \item Instrui usuários a utilizarem sempre a conta institucional \texttt{@academico.ufs.br};
        \item Audita tabelas para assegurar que nenhum registro orfão permaneça no banco.
    \end{itemize}
};
\node[header, above=0.1cm of fala] {O QUE FALA E FAZ?};

% Pain Box: Dores
\node[painbox, below=0.6cm of fala.south west, anchor=north west] (dores) {
    \textbf{\color{danger-red} DORES (FRUSTRAÇÕES E OBSTÁCULOS)}:
    \begin{itemize}\setlength{\itemsep}{1pt}
        \item Demanda repetitiva de atendimento para reset manual de credenciais de acesso;
        \item Cadastros inconsistentes gerados por usuários (ex.: tentar cadastrar e-mails comerciais comuns);
        \item Interfaces sobrecarregadas que não avisam claramente sobre ações destrutivas irreversíveis;
        \item Medo de vazamento acidental de dados pessoais de estudantes universitários.
    \end{itemize}
};

% Gain Box: Ganhos
\node[gainbox, below=0.6cm of fala.south east, anchor=north east] (ganhos) {
    \textbf{\color{hospital-green} GANHOS (DESEJOS E ASPIRAÇÕES)}:
    \begin{itemize}\setlength{\itemsep}{1pt}
        \item Validação estrita e automática de e-mails institucionais (\texttt{@academico.ufs.br});
        \item Segregação arquitetural de matrícula obrigatória exclusiva para discentes;
        \item Redefinição obrigatória de senha no primeiro login delegada ao próprio usuário;
        \item Modais de confirmação em duas etapas para ações destrutivas (exclusão de registros);
        \item Tabelas com paginação universal de 10 itens e ordenação rápida por coluna.
    \end{itemize}
};

\end{tikzpicture}
}
\caption{Mapa de Empatia em Seis Quadrantes: Carlos Eduardo Santos (Administrador)}
\label{fig_mapa_empatia_carlos}
\fonte{Elaborado pelo autor (2026).}
\end{figure}

% ===========================================================================
\section{Síntese Comparativa de Dores e Ganhos}
\label{sec_sintese_mapas}
% ===========================================================================

A Tabela \ref{tab_sintese_mapas_empatia} consolida o cruzamento analítico entre as três personas, discriminando as tensões operacionais e as soluções de design correspondentes implementadas no SIGEA-GTT.

\begin{table}[!ht]
\centering
\caption{Matriz de Síntese Comparativa: Dores, Ganhos e Soluções de Design}
\label{tab_sintese_mapas_empatia}
\footnotesize
\begin{tabularx}{\textwidth}{p{2.6cm} p{3.8cm} p{4.2cm} X}
\toprule
\textbf{Persona} & \textbf{Dores Centrais (Pains)} & \textbf{Ganhos Almejados (Gains)} & \textbf{Solução de Interface no SIGEA-GTT} \\
\midrule
\textbf{Beatriz Matos} \par (Discente) & Ansiedade com cronômetro rígido; sobrecarga de leitura; medo de culpar profissionais. & Auditoria guiada em 20 min; prontuário organizado; foco formativo nas causas sistêmicas. & Cronômetro com alertas visuais discretos aos 15 e 18 min; prontuário em abas estruturadas; ferramentas da qualidade integradas. \\
\addlinespace
\textbf{Prof. Ricardo} \par (Docente) & Exaustão na correção manual de dezenas de papéis; erros de fórmulas no Excel; falta de indicadores. & Espelho automatizado de gabarito; cálculo de taxas do IHI; parecer pedagógico individual. & Comparativo lado a lado de submissão vs. gabarito; cálculo automático de Sensibilidade e $TI_{1000}$; bloqueio de fechamento de turma sem nota. \\
\addlinespace
\textbf{Carlos Santos} \par (Administrador) & Chamados de reset de senha; dados inconsistentes; risco de violação da LGPD. & Validação institucional estrita; integridade referencial; governança sem atrito. & Rejeição automática de e-mails fora de \texttt{@academico.ufs.br}; troca forçada de senha provisória; modais de confirmação em 2 etapas. \\
\bottomrule
\end{tabularx}
\fonte{Elaborado pelo autor (2026).}
\end{table}
